% 控制函数包络：(a) [0,1]上 f_n=x^n（n=1,2,4,8）都在 g=1 之下，积分 1/(n+1)→0；
% (b) f_n=n·1_{(0,1/n)}（n=1..5）积分恒为1；上包络 h=sup_n f_n 在 [1/(n+1),1/n) 上等于 n，且 h>=1/x-1，∫_0^1(1/x-1)dx=∞。
\begin{tikzpicture}[font=\fontsize{8}{10}\selectfont]
\pgfplotsset{every axis/.style={width=55mm,height=42mm,xmin=0,xmax=1.12,
  axis lines=left,axis line style={FigureStroke,line width=.75pt,-{Stealth[length=4pt]}},
  tick style={FigureStroke,line width=.6pt},tick label style={text=FigureInk},
  xlabel={$x$},xlabel style={at={(ticklabel* cs:1)},anchor=north west,text=FigureInk}}}
\begin{axis}[name=a,ymin=0,ymax=1.3,xtick={0,0.5,1},ytick={0,1},title={(a) 共同可积上界$g=1$}]
  \addplot[draw=none,fill=FigureGreenFill] coordinates {(0,0) (0,1) (1,1) (1,0)} \closedcycle;
  \addplot[FigureGreen,line width=1.1pt] coordinates {(0,1) (1,1)};
  \foreach \n in {1,2,4,8}{\edef\temp{\noexpand\addplot[FigureBlue,line width=1.1pt,domain=0:1,samples=80] {x^\n};}\temp}
  \node[text=FigureGreen,anchor=south] at (axis cs:0.5,1) {$g$};
  \node[text=FigureBlue,anchor=west] at (axis cs:0.05,0.85) {$f_n=x^n$};
\end{axis}
\begin{axis}[at={(a.east)},anchor=west,xshift=8mm,ymin=0,ymax=6.4,xtick={0,0.2,0.25,0.333,0.5,1},
  xticklabels={$0$,,,,$0.5$,$1$},ytick={1,2,3,4,5},title={(b) $f_n=n\mathbf 1_{(0,1/n)}$的上包络}]
  \foreach \n in {1,...,5}{\edef\temp{\noexpand\addplot[FigureMuted,line width=.6pt] coordinates {(0,\n) (1/\n,\n) (1/\n,0)};}\temp}
  \addplot[FigureBlue,line width=1.1pt,const plot] coordinates
    {(1/6,5) (1/5,4) (1/4,3) (1/3,2) (1/2,1) (1,1)};
  \addplot[FigureRed,dashed,line width=1.1pt,domain=0.15:1,samples=100] {1/x-1};
  \node[text=FigureBlue,anchor=west] at (axis cs:0.52,1.6) {$h=\sup_nf_n$};
  \node[text=FigureRed,anchor=west] at (axis cs:0.2,5.6) {$1/x-1$};
\end{axis}
\end{tikzpicture}
